% Theme choice:
\usetheme{default}
\usepackage{fontspec}
\usepackage{xcolor}
\usepackage{xeCJK} % 使用 xeCJK 包以支持中文字体
\usepackage{tikz}
\renewcommand{\familydefault}{\sfdefault} % 设置无衬线字体为默认字体
%\usefonttheme{professionalfonts} % 使用 LaTeX 安装的专业字体
%\usepackage{helvet} % 使用 Helvetica 字体
\setbeamertemplate{navigation symbols}{}
% 自定义颜色
\definecolor{customblue}{RGB}{8, 81, 156}

% 设置各部分颜色
\setbeamercolor{title}{fg=customblue}
\setbeamercolor{frametitle}{fg=customblue}
\setbeamercolor{section in toc}{fg=customblue}
\setbeamercolor{subsection in toc}{fg=customblue}
\setbeamercolor{item}{fg=customblue}

% 设置字体大小和样式
\setbeamerfont{title}{series=\bfseries, size=\Large}
\setbeamerfont{frametitle}{series=\bfseries, size=\Large}
\setbeamerfont{author}{size=\scriptsize}
\setbeamerfont{date}{size=\tiny}
\setbeamerfont{subtitle}{size=\footnotesize}
\setbeamerfont{footline}{size=\scriptsize}
\setbeamerfont{block title}{series=\bfseries, size=\small}
\setbeamerfont{block body}{size=\scriptsize}
\usefonttheme{serif} % 使用衬线字体主题
\usepackage{helvet} % 引入 Helvetica 字体
\renewcommand{\familydefault}{\sfdefault} % 将文档的默认字体设为无衬线字体

% 修改 itemize 项的形状为圆点
\setbeamertemplate{itemize item}{\bullet}

% 自定义 section 页面样式
\newcommand{\sectionpagecustom}{
    \begin{center}
        \parbox{0.8\textwidth}{ % 设定整体宽度
            \raggedright\usebeamerfont{section title}\insertsection\par % 左对齐标题
            \vspace{0.2em} % 调整标题和横线之间的距离
            \textcolor{customblue}{\rule{\linewidth}{0.5pt}} % 横线颜色和厚度
        }
    \end{center}
}

% 在目录中为每个 section 标题前添加符号并增加缩进
\setbeamertemplate{section in toc}{\leavevmode\hspace{5em}\llap{\textbullet\hskip0.5em}\inserttocsection}

% 应用自定义 section 页面模板
\setbeamertemplate{section page}{}
\addtobeamertemplate{section page}{\sectionpagecustom}{}

% 自动在每个 section 开始时显示 section 页面
\AtBeginSection{
    \begin{frame}
        \sectionpage
    \end{frame}
}

% 设置页码从目录页开始
\AtBeginDocument{
    % 标题页（不显示页码）
    {
        \setbeamertemplate{footline}{} % 不显示页脚
        \begin{frame}
            \titlepage
        \end{frame}
    }

    % 目录页（开始显示页码，且从1开始计数）
    \setbeamertemplate{footline}{
        \begin{beamercolorbox}[wd=\paperwidth,ht=2.5ex,dp=1ex,leftskip=1em,rightskip=1em]{footline}
            \hfill
            \begin{tikzpicture}
                \draw[fill=customblue, draw=customblue, text=white] circle (0.4cm) node {\tiny \insertframenumber{} / \inserttotalframenumber};
            \end{tikzpicture}
            \hspace{1em}
        \end{beamercolorbox}
    }
    \setcounter{framenumber}{0} % 从目录页开始计数
}
% 自定义页脚，添加当前页码/总页码的圆圈
\setbeamertemplate{footline}{
    \begin{beamercolorbox}[wd=\paperwidth,ht=2.5ex,dp=1ex,leftskip=1em,rightskip=1em]{footline}
        \hfill
        \begin{tikzpicture}
            \draw[fill=customblue, draw=customblue, text=white] circle (0.3cm) node {\tiny \insertframenumber{} / \inserttotalframenumber};
        \end{tikzpicture}
        \hspace{1em}
    \end{beamercolorbox}
}
